\section{An augmented system and analytic estimates}
\label{sec:analytic-estimates}

This section constructs a scalar-Laplacian augmented operator and a
supported estimate that will control differences of transferred physical
displacement--divergence pairs.
Adjoining the displacement divergence to obtain a system with Laplacian
principal part goes back to Weck \cite{Weck1969}; see the account of
Imanuvilov and Yamamoto \cite[Section~1]{ImanuvilovYamamoto2004Carleman}.
The identities below give the formulation used in this proof.
We first work with one of the smoothly extended coefficient pairs from
Lemma~\ref{lem:common-extension}, and suppress its subscript. Set
\[
 m=\mu,\qquad k=\lambda+\mu,\qquad \ell=\lambda+2\mu=k+m.
\]
The energy assumptions imply $m>0$, $k>m/3$, and $\ell>4m/3$.
For an independent vector--scalar pair $U=(u,d)$, define
\begin{equation}
\label{eq:augmented-operators}
\begin{aligned}
 R(u,d)_i
  & =m\Delta u_i+k\partial_i d
   +\sum_{j=1}^3(\partial_jm)(\partial_ju_i+\partial_iu_j)
   +(\partial_i\lambda)d,\\
 S(u,d)
  & =\ell\Delta d+2\nabla k\cdot\nabla d
   +2\nabla m\cdot\Delta u
   +2\sum_{i,j=1}^3(\partial_i\partial_jm)\partial_iu_j
   +(\Delta\lambda)d,\\
 P&=m\Delta+\nabla m\cdot\nabla.
\end{aligned}
\end{equation}
The scalar $d$ will equal $\diverg u$ for physical solutions, but keeping
it independent gives the compatibility identity needed below.

\begin{lemma}
\label{lem:augmented-system}
For every smooth pair $(u,d)$,
\begin{equation}
\label{eq:augmented-identity}
 Lu=R(u,\diverg u),\qquad
 \diverg R(u,d)-S(u,d)=P(\diverg u-d).
\end{equation}
In particular, $Lu=0$ implies $R(u,\diverg u)=S(u,\diverg u)=0$.
The normalized augmented operator
\begin{equation}
\label{eq:normalized-augmented}
 \Lcal U=
 \left(m^{-1}R(u,d),\,
 \ell^{-1}\bigl(S(u,d)-2m^{-1}\nabla m\cdot R(u,d)\bigr)\right)
\end{equation}
has the form
\begin{equation}
\label{eq:laplace-principal-form}
 \Lcal=\Delta\Id_4+\sum_{j=1}^3 A_j(x)\partial_j+A_0(x)
\end{equation}
with smooth matrix coefficients. Consequently, the difference of the
normalized augmented operators for the two coefficient pairs has order
at most one and coefficients supported in $\overline\Omega$.
\end{lemma}

\begin{proof}
The first identity is the coordinate expansion of $\diverg\sigma(u)$.
To verify the second independently of the constraint, write
$v=\diverg u$. Differentiating \eqref{eq:augmented-operators} gives
\[
\begin{aligned}
 \diverg R(u,d)
 ={}&m\Delta v+\nabla m\cdot\nabla v+k\Delta d
       +(\nabla k+\nabla\lambda)\cdot\nabla d\\
    &+2\nabla m\cdot\Delta u
       +2\sum_{i,j=1}^3(\partial_i\partial_jm)\partial_iu_j
       +(\Delta\lambda)d.
\end{aligned}
\]
Subtracting $S$ and using $k=\lambda+m$ gives
$m\Delta(v-d)+\nabla m\cdot\nabla(v-d)$.

For the principal part in \eqref{eq:normalized-augmented}, the only
second derivatives of $u$ in its scalar numerator cancel. More
explicitly, if
\[
 C_m(u)_i=\sum_{j=1}^3(\partial_jm)
                   (\partial_ju_i+\partial_iu_j),
\]
then this numerator is
\[
\begin{aligned}
 S-2m^{-1}\nabla m\cdot R
 ={}&\ell\Delta d+
   (2\nabla k-2km^{-1}\nabla m)\cdot\nabla d\\
 &+2\sum_{i,j=1}^3(\partial_i\partial_jm)\partial_iu_j
   -2m^{-1}\nabla m\cdot C_m(u)\\
 &+\bigl(\Delta\lambda-2m^{-1}\nabla m\cdot\nabla\lambda\bigr)d.
\end{aligned}
\]
This proves \eqref{eq:laplace-principal-form}. The support assertion follows
from equality of the extended coefficients off $\overline\Omega$.
\end{proof}

The order-one difference of the two augmented operators is the forcing
term to be estimated after physical transfer in
Section~\ref{sec:transport-matching}.

Fix a nonzero $\theta\in\C^3$ with $\theta\cdot\theta=0$, and write
\begin{equation}
\label{eq:phase-notation}
 D_\theta=\theta\cdot\nabla,\qquad
 E_\tau(x)=e^{\tau\theta\cdot x},\qquad h=\tau^{-1}.
\end{equation}
For $s\in\R$, the semiclassical Sobolev norm on $\R^3$ is
\[
 \norm{v}_{H_h^s}^2
 =\int_{\R^3}(1+h^2\abs{\xi}^2)^s
                    \abs{\widehat v(\xi)}^2\,\dd\xi.
\]
For vectors, we sum the component norms. These norms and all Hilbert
space adjoints below use the Hermitian structure, whereas products
such as $\theta\cdot\theta$ remain complex bilinear. Constants in the
following estimates may depend on the coefficients, the fixed ball,
and $\theta$, but not on sufficiently large $\tau$.

\begin{proposition}[Supported Carleman estimate]
\label{prop:carleman}
Let $B$ be a bounded ball. For $U\in C_c^\infty(B;\C^4)$ and all
sufficiently small $h>0$,
\begin{equation}
\label{eq:augmented-carleman}
 \norm{E_\tau^{-1}U}_{H_h^1}
 \le Ch\norm{E_\tau^{-1}\Lcal U}_{H_h^{-1}}.
\end{equation}
\end{proposition}

\begin{proof}
Put $\varphi(x)=-\operatorname{Re}(\theta\cdot x)$. A nonzero complex
null vector has nonzero real part, so $\varphi$ is a nonconstant linear
limiting Carleman weight for the Laplacian. For
$\varphi_\varepsilon=\varphi+h\varphi^2/(2\varepsilon)$,
\cite[Lemma~2.1]{SaloTzou2009}, with Sobolev index $s=-1$, gives
\begin{equation}
\label{eq:convexified-carleman}
 \frac{h}{\sqrt\varepsilon}\norm{v}_{H_h^1}
 \le C\norm{e^{\varphi_\varepsilon/h}h^2\Delta
                 e^{-\varphi_\varepsilon/h}v}_{H_h^{-1}},
 \qquad v\in C_c^\infty(B),
\end{equation}
for $h\ll\varepsilon\ll1$, with $C$ independent of sufficiently small
$\varepsilon$. Its statement has a gain of two derivatives and permits
every real $s$. The sign of the Laplacian is immaterial here.

Apply \eqref{eq:convexified-carleman} componentwise. A first-order term
in \eqref{eq:laplace-principal-form} contributes
\[
 e^{\varphi_\varepsilon/h}h^2A_j\partial_j
       e^{-\varphi_\varepsilon/h}v
 =hA_j(h\partial_jv)-hA_j(\partial_j\varphi_\varepsilon)v.
\]
Its $H_h^{-1}$ norm is bounded by $C_1h\norm{v}_{H_h^1}$,
uniformly when $h/\varepsilon$ is small. The zeroth-order terms have
the same bound. Choose $\varepsilon>0$ so small that these terms are
absorbed by the left side of \eqref{eq:convexified-carleman}, and then
choose $h$ sufficiently small relative to this fixed $\varepsilon$.
We obtain
\[
 \norm{e^{\varphi_\varepsilon/h}U}_{H_h^1}
 \le C_\varepsilon h
       \norm{e^{\varphi_\varepsilon/h}\Lcal U}_{H_h^{-1}}.
\]

The factor $e^{\varphi^2/(2\varepsilon)}$ is now fixed. It and its
inverse, smoothly localized around $\overline B$, are bounded
multipliers on $H_h^1$, uniformly for $0<h\le1$, by the product rule.
They are bounded on $H_h^{-1}$ by duality. Removing this factor gives
the same estimate with $\varphi$ in place of $\varphi_\varepsilon$.
The multiplier constants may depend on $\varepsilon$; no subsequent
absorption involves them.

Finally, write $\beta=\operatorname{Im}\theta$.
Multiplication by $e^{-i\beta\cdot x/h}$ shifts the semiclassical
frequency $h\xi$ by the fixed vector $-\beta$. The weights
$1+\abs{h\xi}^2$ and $1+\abs{h\xi-\beta}^2$ are comparable,
so this multiplier and its inverse are uniformly bounded on
$H_h^{\pm1}$. Since
$E_\tau^{-1}=e^{\varphi/h}e^{-i\beta\cdot x/h}$, this proves
\eqref{eq:augmented-carleman}.
\end{proof}

\begin{corollary}
\label{cor:physical-solvability}
For the same $B$ and $\theta$, and all sufficiently large $\tau$,
every $w\in C_c^\infty(B;\C^3)$
satisfies
\begin{equation}
\label{eq:physical-test-estimate}
 \norm{E_\tau^{-1}w}_{L^2(B)}
 \le C\norm{E_\tau^{-1}Lw}_{L^2(B)}.
\end{equation}
Moreover, for every $f\in L^2(B;\C^3)$ there exists
$z\in L^2(B;\C^3)$ satisfying, distributionally,
\begin{equation}
\label{eq:physical-solvability}
 E_\tau^{-1}L(E_\tau z)=f\quad\text{in }B,
 \qquad \norm{z}_{L^2(B)}\le C\norm{f}_{L^2(B)}.
\end{equation}
\end{corollary}

\begin{proof}
Apply Proposition~\ref{prop:carleman} to $(w,\diverg w)$.
By \eqref{eq:augmented-identity}, its augmented forcing is formed
from $Lw$ and $\diverg(Lw)$. If $F=E_\tau^{-1}Lw$, then
\[
 E_\tau^{-1}\diverg(Lw)=\diverg F+\tau\theta\cdot F.
\]
The Fourier definition of the norms gives
$\norm{\partial_jF}_{H_h^{-1}}\le h^{-1}\norm{F}_{L^2}$.
Smooth coefficient multipliers, localized near $\overline B$, are
uniformly bounded on $H_h^{-1}$. Hence
\[
 \norm{E_\tau^{-1}\Lcal(w,\diverg w)}_{H_h^{-1}}
 \le Ch^{-1}\norm{F}_{L^2}.
\]
Taking the displacement component on the left side of
\eqref{eq:augmented-carleman} proves
\eqref{eq:physical-test-estimate}.

Let $A_\theta=E_\tau^{-1}LE_\tau$ on test functions in $B$.
Substituting $w=E_\tau v$ in
\eqref{eq:physical-test-estimate} gives
$\norm{v}_{L^2}\le C\norm{A_\theta v}_{L^2}$.
Formal self-adjointness of $L$ gives
\[
 A_\theta^*
 =e^{\tau\overline\theta\cdot x}L
                      e^{-\tau\overline\theta\cdot x}
 =A_{-\overline\theta}.
\]
The same estimate for the null direction $-\overline\theta$ thus
holds for $A_\theta^*$. Use the $L^2$ inner product linear in its
first argument. The functional
\[
 A_\theta^*v\longmapsto (v,f)_{L^2(B)},
 \qquad v\in C_c^\infty(B;\C^3),
\]
is well-defined and has norm at most $C\norm{f}_{L^2}$. Extend it by
the Hahn--Banach Theorem and represent it as $(\,\cdot\,,z)_{L^2}$.
Then $(A_\theta^*v,z)=(v,f)$ for every test function $v$, which is
the distributional equation in \eqref{eq:physical-solvability}.
The norm of the representing vector gives the asserted bound.
\end{proof}

\begin{lemma}[Scaled interior estimate]
\label{lem:scaled-interior}
Let $B'\Subset B$. There exists $h_{B'}>0$ such that, whenever
$0<h=\tau^{-1}<h_{B'}$ and $z,f\in L^2(B;\C^3)$ satisfy
$E_\tau^{-1}L(E_\tau z)=f$ distributionally in $B$,
\begin{equation}
\label{eq:scaled-interior}
 \sum_{j=0}^2 h^j\norm{\nabla^jz}_{L^2(B')}
 \le C_{B'}\bigl(\norm{z}_{L^2(B)}
                         +h^2\norm{f}_{L^2(B)}\bigr).
\end{equation}
If $f$ is smooth, then $z$ is smooth in $B$ for each fixed $\tau$.
\end{lemma}

\begin{proof}
The positive principal coefficient form of $L$ is
\[
 m\abs{\xi}^2\abs{v}^2+k\abs{\xi\cdot v}^2,
 \qquad \xi\in\R^3,\quad v\in\C^3.
\]
Thus its ellipticity is uniform on $\overline B$. Conjugation leaves
the principal part unchanged and introduces coefficients of size
$O(h^{-1})$ and $O(h^{-2})$ in orders one and zero. On a ball
$B(x_0,2h)\subset B$, set $x=x_0+hy$ and multiply the equation by
$h^2$. The resulting operator on the fixed ball $B(0,2)$ has a
uniformly strongly elliptic principal part, while its coefficients
and their derivatives in $y$ are uniformly bounded. As
$(x_0,h)$ ranges over $\overline B'\times[0,h_0]$, these operators
form a compact smooth coefficient family with a common ellipticity
bound. The local parametrix estimate persists under a small
coefficient perturbation, so a finite cover gives a uniform constant.
Interior $H^2$ regularity for distributional $L^2$ solutions of smooth
elliptic systems is given by
\cite[Theorem~15.1 and Remark~15.2]{Dyatlov2026}. The same parametrix
argument with matrix symbols gives the localized estimate for systems
corresponding to the scalar Proposition~15.6 of that source. With the
uniform constants just obtained, it yields
\[
 \sum_{j=0}^2 h^j\norm{\nabla^jz}_{L^2(B(x_0,h))}
 \le C\bigl(\norm{z}_{L^2(B(x_0,2h))}
                      +h^2\norm{f}_{L^2(B(x_0,2h))}\bigr).
\]
For small $h$, cover $B'$ by such inner balls with uniformly bounded
overlap of the doubled balls. Squaring and summing proves
\eqref{eq:scaled-interior}. Iterated interior elliptic regularity,
with $\tau$ fixed, proves the final assertion.
\end{proof}
